\documentclass[10pt,a4paper]{article}
\usepackage{amsmath,amssymb,amsthm,graphicx}
\newif\iflipics\lipicsfalse
\usepackage{amsmath}
\usepackage{amsfonts}
\usepackage{amssymb}
\usepackage{amsthm}
\usepackage{xspace}
\usepackage[dvipsnames]{xcolor}
\usepackage{tabularx}
\usepackage{ragged2e}
\usepackage{nicefrac}
\usepackage{bbm}
\usepackage{caption}
\usepackage{graphicx}
\usepackage{enumitem}
\usepackage{soul}
\usepackage{hyperref}
\definecolor{darkgreen}{rgb}{0,0.4,0}
\usepackage{algorithm}
\usepackage[commentColor=darkgreen]{algpseudocodex}
\newcommand{\LineComment}{\LComment}
% \algrenewcomment[1]{\hfill \textit{\textcolor{darkgreen}{$\triangleright$ #1}}}
% \algnewcommand{\LineComment}[1]{\State \textit{\textcolor{darkgreen}{$\triangleright$ #1}}}
\algnewcommand{\IfThenElse}[3]{% \IfThenElse{<if>}{<then>}{<else>}
  \State \algorithmicif\ #1\ \algorithmicthen\ #2\ \algorithmicelse\ #3}
\algdef{SE}[SUBALG]{Indent}{EndIndent}{}{\algorithmicend\ }%
\makeatletter
\renewcommand{\ALG@name}{Listing}
\makeatother
\usepackage{cleveref}
\crefname{listing}{listing}{listings}
\crefname{equation}{eq.}{eqs.}

% must be after cleveref
\newtheorem{theorem}{Theorem}[section]
\newtheorem{lemma}[theorem]{Lemma}
\newtheorem{corollary}[theorem]{Corollary}
\newtheorem{fact}{Fact}[section]
\newtheorem{conj}[theorem]{Conjecture}
\newtheorem{question}[theorem]{Question}
\newtheorem{claim}[theorem]{Claim}
\newtheorem{obs}[theorem]{Observation}
\newtheorem{proposition}[theorem]{Proposition}
\theoremstyle{definition}
\newtheorem{definition}[theorem]{Definition}
\newtheorem{example}{Example}
\newtheorem{results}{Result Family}
\newtheorem{result}{Result}
\theoremstyle{remark}
\newtheorem*{remark}{Remark}
\makeatletter
\newcommand\squeezepar{\@startsection{paragraph}{4}%
{\z@}{1.5ex \@plus1ex \@minus.2ex}{-1em}{\normalfont\normalsize\bfseries}}
\newcommand{\mypara}[1]{\squeezepar{{#1}.}}
\renewcommand{\paragraph}[1]{\squeezepar{{#1}.}}
\makeatother
\graphicspath{{items/}{./}}
\newcommand{\EE}{\mathbb{E}}
\newcommand{\FF}{\mathbb{F}}
\newcommand{\NN}{\mathbb{N}}
\newcommand{\RR}{\mathbb{R}}
\newcommand{\ZZ}{\mathbb{Z}}
\newcommand{\indic}{\mathbbm{1}}
\newcommand{\cA}{\mathcal{A}}
\newcommand{\cB}{\mathcal{B}}
\newcommand{\cC}{\mathcal{C}}
\newcommand{\cD}{\mathcal{D}}
\newcommand{\cE}{\mathcal{E}}
\newcommand{\cF}{\mathcal{F}}
\newcommand{\cG}{\mathcal{G}}
\newcommand{\cH}{\mathcal{H}}
\newcommand{\cI}{\mathcal{I}}
\newcommand{\cJ}{\mathcal{J}}
\newcommand{\cM}{\mathcal{M}}
\newcommand{\cO}{\mathcal{O}}
\newcommand{\cP}{\mathcal{P}}
\newcommand{\cQ}{\mathcal{Q}}
\newcommand{\cR}{\mathcal{R}}
\newcommand{\cS}{\mathcal{S}}
\newcommand{\cU}{\mathcal{U}}
\newcommand{\cV}{\mathcal{V}}
\newcommand{\cW}{\mathcal{W}}
\newcommand{\cX}{\mathcal{X}}
\newcommand{\cY}{\mathcal{Y}}
\newcommand{\cZ}{\mathcal{Z}}

\newcommand{\tG}{\widetilde{G}}
\newcommand{\tS}{\widetilde{S}}
\newcommand{\tO}{\widetilde{O}}
\newcommand{\tP}{\widetilde{P}}
\newcommand{\tQ}{\widetilde{Q}}
\newcommand{\tf}{\widetilde{f}}
\newcommand{\tTheta}{\widetilde{\Theta}}
\newcommand{\tOmega}{\widetilde{\Omega}}
\newcommand{\tM}{\widetilde{M}}
\newcommand{\tmu}{\widetilde{\mu}}

\newcommand{\hw}{\hat{w}}
\newcommand{\hb}{\hat{b}}
\newcommand{\hP}{\hat{P}}
\newcommand{\hf}{\hat{f}}
\newcommand{\hQ}{\hat{Q}}

\newcommand{\ceil}[1]{{\left\lceil{#1}\right\rceil}}
\newcommand{\floor}[1]{{\left\lfloor{#1}\right\rfloor}}
\newcommand{\norm}[1]{{\left\Vert{#1}\right\Vert}}
\DeclareSymbolFont{newfont}{OML}{cmm}{m}{it}% Computer Modern math font
\DeclareMathSymbol{\emptystream}{3}{newfont}{15}% Symbol 15 "epsilon"
\let\epsilon\varepsilon
\renewcommand{\emptyset}{\varnothing}

\newcommand{\avoid}{\textsc{avoid}\xspace}
\newcommand{\mif}{\textsc{mif}\xspace}
\newcommand{\miflong}{\textsc{MissingItemFinding}\xspace}
\newcommand{\fco}{\textsc{fco}\xspace}
\newcommand{\fao}{\textsc{fao}\xspace}
\newcommand{\init}{\mathrm{init}\xspace}
\newcommand{\poly}{\mathrm{poly}\xspace}
\newcommand{\polylog}{\mathrm{polylog}\xspace}
\newcommand{\Algs}{\textsc{Algs}\xspace}

\newcommand{\repeatsc}{\textsc{repeat}\xspace}
\newcommand{\unsafesc}{\textsc{unsafe}\xspace}
\newcommand{\incompletesc}{\textsc{incomplete}\xspace}
\newcommand{\failsc}{\textsc{fail}\xspace}
\newcommand{\emptysc}{\textsc{empty}\xspace}
\newcommand{\timeoutsc}{\textsc{timeout}\xspace}
\newcommand{\bigsc}{\textsc{big}\xspace}
\newcommand{\abort}{\textsc{abort}\xspace}
\newcommand{\sort}{\textsc{sort}\xspace}
\newcommand{\abortsc}{\textsc{abort}\xspace}
\newcommand{\errorsc}{\textsc{error}\xspace}

\newcommand{\rl}{\mathbf{\ell}}             % stream length
\newcommand{\stream}{\tau}                  % generic input stream
\newcommand{\morestream}{\varphi}           % ``extension'' of \stream
\newcommand{\concat}{\mathbin{\circ}}       % concatenation of sequences
\newcommand{\trans}{T}                      % transition function of automaton
\newcommand{\SSD}[2]{\textsc{seqs}(#1,#2)}  % sorted sequences of distinct things
\newcommand{\Outs}{\textsc{Outs}\xspace}    % output seq given start state + adversary
\newcommand{\Mstk}{\Delta}                  % min max probability of mistake

\newcommand{\todo}[1]{%
  \marginpar{\color{red}!!!}[{\color{red}\bf  #1}]}
\newcommand{\gobble}[1]{}
\crefname{originalsection}{original section}{original sections}
\Crefname{originalsection}{Original section}{Original sections}
\makeatletter
\newlabel{subsec:pd-lb-overview}{{2.4}{source}{}{sourcepointer-1}{}}
\newlabel{subsec:pd-lb-overview@cref}{{[originalsection][0][]2.4}{[1][0][]source}{}{}{}}
\makeatother
\crefname{originalsection}{original section}{original sections}
\Crefname{originalsection}{Original section}{Original sections}
\makeatletter
\newlabel{subsec:rt-lb-overview}{{2.2}{source}{}{sourcepointer-2}{}}
\newlabel{subsec:rt-lb-overview@cref}{{[originalsection][0][]2.2}{[1][0][]source}{}{}{}}
\makeatother
\title{Space lower bounds for pseudo-deterministic missing-item algorithms even with a random oracle}
\author{Amit Chakrabarti, Manuel Stoeckl}
\date{doi:10.4230/LIPIcs.CCC.2024.28; CC BY-SA4.0}
\begin{document}
\maketitle
\noindent\textit{Editorial note.} Original source: Amit Chakrabarti and Manuel Stoeckl, arXiv:2310.03634v2 (2024-07-03), https://arxiv.org/abs/2310.03634v2. License: Creative Commons Attribution-ShareAlike4.0, https://creativecommons.org/licenses/by-sa/4.0/.\par\medskip
\noindent\textit{Editorial note.} Changes: selected complete author theorem/proof/context excerpts and original diagram/macros/bibliography in a standalone layout; omitted unrelated sections and author frontmatter. No mathematical corrections or additional proofs. This adapted presentation retains CC BY-SA4.0.\par\medskip
\noindent\textit{Allowed standard prerequisites (editorial).}\begin{itemize}
\item One-way avoid communication lower bound. Author Theorem lem:avoid-lb explicitly cites ChakrabartiGS22; exact problem and statement retained.
\end{itemize}\medskip
\ifdefined\hypertarget\hypertarget{sourcepointer-1}{}\fi
\noindent\textit{Source reference.} Original section 2.4: number-theory/sources/ccc-2024-28/native/main-content.tex, source lines 758--758. Optional introductory overview; all canonical-output algorithm and proof steps are included.\par\medskip
\ifdefined\hypertarget\hypertarget{sourcepointer-2}{}\fi
\noindent\textit{Source reference.} Original section 2.2: number-theory/sources/ccc-2024-28/native/main-content.tex, source lines 615--615. Optional reference used only when introducing the avoid problem; the exact problem and cited bound are retained.\par\medskip
In this paper, we prove that for certain parameter regimes, \miflong
(henceforth, $\mif$) is such a problem. In the problem $\mif(n,\rl)$, the
input is a stream $\langle e_1,\ldots,e_\rl \rangle$ of $\rl$ integers, not
necessarily distinct, with each $e_i \in \{1,\ldots,n\}$, where $1 \le \rl \le
n$. The goal is as follows: having received the $i$th integer, output a number
$v$ in $\{1,\ldots,n\} \setminus \{e_1,\ldots,e_i\}$.  We will be mostly interested in the setting $\rl = o(n)$, so
the ``trivial'' upper bound on the space complexity of $\mif(n,\rl)$ is
$O(\rl \log n)$, achieved by the deterministic algorithm that simply stores
the input stream as is.

\subsection{Groundwork for Our Results}\label{subsec:models}

To state our results about $\mif$, we need to introduce some key terminology.
Notice that \mif is a {\em tracking problem}: an output is required after
reading each input.%
%
\footnote{We do not consider algorithms with a ``one-shot'' guarantee, to only
be correct at the end of the stream, because a) the adversarial setting
requires tracking output b) for \mif and most other problems the difference
in space complexity is generally small.} 
%
Thus, we view streaming algorithms as generalizations of finite state
(Moore-type) machines. An algorithm $\cA$ has a finite set of states $\Sigma$
(leading to a space cost of $\log_2|\Sigma|$), a finite input set $\cI$, and a
finite output set $\cO$.  It has a transition function $\trans \colon \Sigma
\times \cI \times \cR \to \Sigma$ indicating the state to switch to after
receiving an input, plus an output function $\gamma \colon \Sigma \times \cR
\to \cO$ indicating the output produced upon reaching a state. How the final
parameter (in $\cR$) of $\trans$ and $\gamma$ is used depends on the type of
randomness. We consider four cases, leading to four different models of
streaming computation.

\iflipics\begin{itemize}\else
\begin{itemize}[left=2ex]\fi
  %
  \item \emph{Deterministic.~} The initial state of the algorithm is a fixed
  element of $\Sigma$, and $\trans$ and $\gamma$ are deterministic (they do not
  depend on the parameter in $\cR$). 
  %
  \item \emph{Random seed.~} The initial state is drawn from a distribution
  $\cD$ over $\Sigma$, and $\trans$ and $\gamma$ are deterministic. This models
  the situation that all random bits used count towards the algorithm's space
  cost.
  %
  \item \emph{Random tape.~} The initial state is drawn from a distribution
  $\cD$ over $\Sigma$.\footnote{Requiring that this model use a fixed initial
  state could make some algorithms use one additional ``\textsc{init}''
  state.} The space $\cR$ is a sample space; when the algorithm receives an
  input $e \in \cI$ and is at state $\sigma \in \Sigma$, it chooses a random
  $\rho \in \cR$ independent of all previous choices and moves to state
  $\trans(e, \sigma, \rho)$. However, $\gamma$ is deterministic.%
  %
  \footnote{Alternatively, we could associate a \emph{distribution} of outputs
  to each state, or a function mapping (input, state) pairs to outputs. As
  these formulations are slightly more complicated to prove things with, and
  only affect the space usage of \miflong algorithms by an additive $O(\log n
  + \log \frac{1}{\delta})$ amount, we stick with the one state = one output
  convention.} %
  %
  This models the situation that the algorithm can make random decisions at
  any time, but it cannot remember past random decisions without recording
  them (which would add to its space cost). 
  %
  \item \emph{Random oracle.~} The initial state is fixed; $\cR$ is a
  sample space. A specific $R \in \cR$ is drawn at the start of the algorithm
  and stays the same over its lifetime. When the algorithm is at state
  $\sigma$ and receives input $e$, its next state is $\trans(e, \sigma, R)$.
  The output given at state $\sigma$ is $\gamma(\sigma, R)$.  This models the
  situation that random bits are essentially ``free'' to the algorithm; it can
  read from a long random string which doesn't count toward its space cost and
  which remains consistent over its lifetime. A random oracle algorithm can be
  interpreted as choosing a random deterministic algorithm, indexed by $R$,
  from some family.  
  %
\end{itemize}

These models form a rough hierarchy; they have been presented in (almost)
increasing order of power. Every $z$-bit ($2^z$-state) deterministic algorithm
can be implemented in any of the random models using $z$ bits of space; the
same holds for any $z$-bit random seed algorithm.  Every $z$-bit random tape
algorithm has a corresponding $(z + \log\rl)$-bit random oracle
algorithm---the added space cost is because for a random oracle algorithm to
emulate a random tape algorithm, it must have a way to get ``fresh''
randomness on each turn.%
%
\footnote{An alternative, which lets one express $z$-bit random tape
algorithms using a $z$-bit random oracle variant, is to assume the random
oracle algorithm has access to a clock or knows the position in the stream for free;
both are reasonable assumptions in practice.}

Streaming algorithms are also classified by the kind of correctness guarantee
they provide. Recall that we focus on ``tracking''
algorithms~\cite{BenEliezerJWY20}; they present an output after reading each
input item and this {\em entire sequence} of outputs must be correct.  Here
are three possible meanings of the statement ``algorithm $\cA$ is
$\delta$-error'' (we assume that $\cA$ handles streams of length $\rl$ with
elements in $\cI$ and has outputs in $\cO$):

\iflipics\begin{itemize}\else
\begin{itemize}[left=2ex]\fi
  %
  \item \emph{Static setting.~} For all inputs $\stream \in \cI^\rl$,
  running $\cA$ on $\stream$ produces incorrect output with probability $\le
  \delta$.
  %
  \item \emph{Adversarial setting.~} For all (computationally unbounded)
  adaptive adversaries $\alpha$ (i.e., for all functions
  $\alpha \colon \cO^\star \to \cI$),\footnote{By the minimax
  theorem, it suffices to consider deterministic adversaries.} running $\cA$
  against $\alpha$ will produce incorrect output with probability $\le
  \delta$.
  %
  \item \emph{Pseudo-deterministic setting.~} There exists a canonical output
  function $f \colon \cI^\star \rightarrow \cO$ producing all correct outputs
  so that, for each $\stream \in \cI^\rl$, $\cA(\stream)$ fails to output
  $f(\stream)$ with probability $\le \delta$.
  %
\end{itemize}

Algorithms for the static setting are called ``classic'' streaming algorithms;
ones for the adversarial setting are called ``adversarially robust'' streaming
algorithms. All pseudo-deterministic algorithms are adversarially robust, and
all adversarially robust algorithms are also classic.

As a consequence of Newman's theorem~\cite{Newman91}, any random oracle or
random tape algorithm in the static setting with error $\delta$ can be
emulated using a random seed algorithm with only $\epsilon$ increase in error
and an additional $O(\log \rl + \log \log |\cI| + \log \frac{1}{\epsilon
\delta})$ bits of space. However, the resulting algorithm is non-constructive.


\section{Preliminaries}\label{sec:prelim}

\mypara{Notation} Throughout this paper, $\log x = \log_2 x$, while $\ln x =
\log_e x$. The set $\NN$ consists of all positive integers; $[k] :=
\{1,2,\ldots,k\}$; and $[a,b)$ is a half open interval of real numbers. For a
condition or event $E$, the symbol $\indic_{E}$ takes the value $1$ if $E$
occurs and $0$ otherwise. The sequence (stream) obtained by concatenating
sequences $a$ and $b$, in that order, is denoted $a \concat b$. For a set $S$
of elements in a totally ordered universe, $\sort(S)$ denotes the sequence of
elements of $S$ in increasing order; $\binom{S}{k}$ is the set of $k$-element
subsets of $S$; and $\SSD{S}{k} = \{\sort(Y) : Y \in \binom{S}{k}\}$. We
sometimes extend set-theoretic notation to vectors and sequences; e.g., for $y
\in [n]^t$, write $y \subseteq S$ to mean that $\forall i \in [t]: y_i \in S$.
For a set $X$, $\triangle[X]$ denotes the set of probability distributions
over $X$, while $A \in_R X$ indicates that $A$ is chosen uniformly at random
from $X$. When naming probability distributions, we will use either
calligraphic letters (e.g., $\cA,\cD$) or the letter $\mu$; $A \sim \mu$ means
that $A$ is drawn from the distribution $\mu$.

As outlined in \Cref{subsec:rt-lb-overview}, $\avoid(m,a,b)$ is a one-way
communication problem, wherein player Alice has a set $A \in
\smash{\binom{[m]}{a}}$, and should send a short message to player Bob, who
should use the message to output a set $B \in \smash{\binom{[m]}{b}}$ that is
disjoint from $A$. In the randomized $\delta$-error setting, this disjointness
should hold with probability $\ge 1-\delta$.

\begin{theorem}[\avoid communication lower bound, from \cite{ChakrabartiGS22}]\label{lem:avoid-lb}
  Suppose there exists a randomized protocol for $\avoid(m,a,b)$, in which
  Alice communicates $\le K$ bits, that is $\delta$-error either on a
  worst-case input or when Alice's input is chosen uniformly at random from
  $\binom{[t]}{a}$. Then
  \begin{align*}
    K \ge \frac{a b}{t \ln 2} + \log(1-\delta) \,.
  \end{align*}
\end{theorem}

\begin{lemma}[Error amplification by majority vote]\label{lem:error-reduction-by-vote}
  Let $\epsilon \le \delta \le 1/3$. Say $X$ is a random variable, and $v$ a value with $\Pr[X = v] \ge 1 - \delta$. If $X_1,\ldots,X_p$ are independent copies of $X$, then the most common value in $(X_1,\ldots,X_p)$ will be $v$ with probability $\ge 1 - \epsilon$, for $\epsilon = \left(2\delta\right)^{p / 30}$.
\end{lemma}

\noindent The above is a standard lemma, useful for trading error for space
for algorithms with a single valid output.

\begin{proof}[Proof of \Cref{lem:error-reduction-by-vote}]
  For each $i \in [p]$, let $Y_i$ be the random indicator variable for the event that $X_i \ne v$. Let $\alpha = \frac{1}{2 \delta} - 1$.The probability that $v$ is not the most common element can be bounded by the probability that it is the not the majority element; by a Chernoff bound, this is:
  \begin{align*}
    \Pr[\sum_{i\in[p]} Y_i \ge \frac{1}{2} p] &= \Pr[\sum_{i\in[p]} Y_i \ge (1 + \alpha) \delta p] \le \exp\left( - ((1+\alpha) \ln(1+\alpha) - \alpha) \delta p\right) \\
      &\le \exp\left( - 0.073 ((1+\alpha) \ln(1+\alpha) ) \delta p\right) \qquad \text{since $\alpha \ge 1/6$} \\
      &\le \exp\left( - \frac{0.073 }{2\delta} \ln\frac{1}{2\delta} \delta p\right) = \left(2\delta\right)^{0.036 p} \,. \qedhere
  \end{align*}
\end{proof}

\section{The Pseudo-Deterministic and Random Seed Lower Bounds}\label{sec:pd-lb}

In this section, we prove a space lower bound for pseudo-deterministic
streaming algorithms; in particular, for the most general (random oracle) type
of them. See \Cref{subsec:pd-lb-overview} for a high-level plan of the proof.

Let $\cA$ be a random-oracle pseudo-deterministic algorithm for $\mif(n, \rl)$
using $z$ bits of state, which has worst case failure probability $\delta \le
\frac{1}{3}$. Let $\Pi \colon [n]^\star \to [n]$ be the function giving the
{\em canonical output} of $\cA$ after processing a stream (as was defined in
\Cref{subsec:models}), and let $S = \Pi([n]^\rl)$ be the set of all canonical
outputs at time $\rl$ (for this proof we can ignore outputs at times $<\rl$).
Clearly $|S| \le n$, and we will also prove (see \Cref{lem:pd-s-le-2zp1}) that
$|S| \le 2^{z+1}$. It is easy to see that $|S| \ge \rl + 1$ and that, since
$\mif(n,\rl)$ is nontrivial, correct algorithms must have $z \ge 1$.

The main proof in this section only applies to algorithms with very low error
(potentially as small as $1/n^{\Omega(\log n)}$). To ensure that we are
working with an algorithm with error this small, we will first apply
\Cref{lem:error-reduction-by-vote}, using $p$ independent instances of $\cA$,
where $p \ge 1$ is an integer chosen later. This will produce an $\epsilon$-error
algorithm $\cB$ that uses $z p$ bits of space, where $\epsilon \le
(2\delta)^{p/30}$. Moreover, the canonical outputs of $\cB$ will still be
given by $\Pi$.

We may view $\cB$ as a distribution over deterministic streaming algorithms;
each such algorithm, obtained by fixing $\cB$'s oracle random string, produces
outputs according to a function of the form $A \colon [n]^\star \to [n]$.
Since $\cB$ is pseudo-deterministic, we have the following property:
\begin{align}
  \forall x \in [n]^{\le \rl}: \Pr_{A \sim \cB} [A(x) \ne \Pi(x)] \le \epsilon \,. 
  \label{eq:pd-fn-prop}
\end{align}

It is time to give the details of the key procedure \textsc{FindCommonOutputs}
(abbreviated as $\fco$), outlined earlier in \Cref{subsec:pd-lb-overview},
that underpins our proof.  The input stream positions, from $1$ to $\rl$, are
split into $d = \Theta(\rl / (zp))$ consecutive intervals, of lengths $t_d,
t_{d-1}, \ldots, t_1$, in order, where $t_d = t_{d-1} = \cdots = t_2 =
\Theta(zp)$ and $t_1 = \rl - \sum_{k=2}^{d} t_k$; the constants are chosen
such that $t_1 \ge \rl/2$. A call to \fco takes the form $\fco(B, C, x,
k)$, where $B \colon [n]^\rl \to [n]$ is an output function, $C$ is a random
collection (matrix) of thresholds, $k \in [d]$, and $x \in [n]^{t_d + \cdots +
t_{k+1}}$ is a stream prefix. (Notice that when $k = d$, we must have $x =
\emptystream$.) By design \fco is recursive, bottoming out at $k = 1$, and
always returns a subset of $S$ of cardinality $w_k$, where
\begin{align}
  w_k := 2^{k-1} (t_1 + 1) \,. \label{eq:wk-def}
\end{align}
The precise logic of \fco is given in \Cref{alg:pd-set-finding}. Throughout
this section, we use $\SSD{Z}{t}$ to denote the set of all length-$t$ {\em
sorted sequences} of {\em distinct} elements of $Z$. Note that
\[
  |\SSD{Z}{t}| = \binom{|Z|}{t} \,.
\]

\begin{algorithm}[htb]
  \caption{The procedure to compute a set for \Cref{lem:pd-main-recursion}}\label[listing]{alg:pd-set-finding}

  \begin{algorithmic}[1]
    \Statex Let $t_1,\ldots,t_d$, $w_1,\ldots,w_d$ be integer parameters, and $S$ the set of valid outputs
    \Statex
    \Statex \underline{\textsc{FindCommonOutputs($B$, $C$, $x$, $k$)}} \Comment{abbreviated as $\fco(B,C,x,k)$}
    \LineComment{Inputs: function $B \colon [n]^\rl \to [n]$, 
      matrix $C \in [1,2)^{d \times \NN}$, stream prefix $x \in [n]^{t_k + \cdots + t_1}$}
    \LineComment{Output: a subset of $S$ of size $w_k$}
    
    \If{$k = 1$} 
        \State $e_0 \gets B(x \concat \langle 1,1,\ldots,1\rangle)$\label{step:pd-base-case-start}
        \For{$i$ in $1,\ldots,t_1$} 
          \State $e_i \gets B(x \concat \langle e_0,\ldots,e_{i-1},1,\ldots,1\rangle)$ \label{step:pd-base-case-mid}
        \EndFor
        \If{$e_0, \ldots,e_{t_1}$ are all distinct}
          \State \textbf{return} $\{e_0,e_1,\ldots,e_{t_1}\}$
          \Comment{identify $w_1$ distinct possible outputs}
        \EndIf
        \State \textbf{return} arbitrary subset of $S$ of size $w_1$ (failure)
    \Else
        \For{each $y \in \SSD{[n]}{t_k}$}
            \State $T_{y} \gets \textsc{FindCommonOutputs}(B, C, x \concat y, k-1)$ \label{step:pd-compute-t}
            \Comment{note $|T_{y}| = w_{k-1}$}
        \EndFor
    
        \State $Q_0 \gets T_{\langle 1,2,\ldots,t_k \rangle}$
        \For{$h$ in $1,2,3,4$}\label{step:pd-for-loop-start}
            \LineComment{gather statistics and find common elements among the sets $T_{y}$}
            \For{each $j\in S$}
                \State $f_{j}^{(h)} \gets \left| \left\{ y \in \SSD{Q_{h-1}}{t_k} :
                  j \in T_{y} \right\} \right|$ \label{step:pd-def-f}
                  \Comment{count frequencies}
            \EndFor
            \State $\theta \gets C_{k,h} w_{k-1} / (16 |S|)$ \Comment{set random threshold}
            \State $P_h \gets \left\{j \in S : f_{j}^{(h)} \ge 
              \theta \binom{|Q_{h-1}|}{t_k} \right\}$ \label{step:pd-comp-ph}
              \Comment{identify ``sufficiently common'' elements}
            \State $Q_h \gets Q_{h-1} \cup P_h$
            \If{$|Q_h| \ge w_k$}
              \State \textbf{return} the $w_k$ smallest elements in $Q_h$ \label{step:pd-good-exit}
            \EndIf
        \EndFor \label{step:pd-for-loop-end}
        \State \textbf{return} arbitrary subset of $S$ of size $w_k$ (failure)\label{step:pd-bad-exit}
    \EndIf
  \end{algorithmic}
  
\end{algorithm}

A key property of \fco is that for all $x \in [n]^{t_d + \cdots + t_{k+1}}$,
if $C$ is chosen uniformly at random from $[1,2)^{d \times \NN}$, and $A \sim
\cB$, then w.h.p. $\fco(A, C, x, k)$ produces the same set as $\fco(\Pi, C, x,
k)$. Another important property is that $\fco(\Pi, C, x, k)$ always produces
$w_k$ distinct elements in $S$ without a failure (i.e., returns using
\Cref{step:pd-good-exit}, not \Cref{step:pd-bad-exit}). These properties
are formally proved in \Cref{lem:pd-main-recursion}, to follow.  In
particular, $\fco(\Pi, C, \emptystream, d)$ produces $w_d$ distinct elements
in $S$, showing that $|S| \ge w_d$. Since $w_d$ is a function of $z$,
combining this with upper bounds on $|S|$ lets us solve for a lower bound on
$z$.

We elaborate on how we achieve the two key properties of $\fco$. First, to
ensure $\fco(\Pi,C,x,k)$ equals $\fco(A,C,x,k)$ for random $A \sim \cB$ and $C
\in_R [1,2)^{d \times \NN}$, we use a standard random threshold trick when
computing the set $P_h$ on \Cref{step:pd-comp-ph}. The recursive calls to
$\fco(A, C, x \concat y, k-1)$ on \Cref{step:pd-comp-ph} do not always
match the outputs of $\fco(\Pi, C, x \concat y, k-1)$; as a result, the
element frequency vector $f^{(h)}$ from \Cref{step:pd-def-f} may have
random noise when computed using $A$ instead of using $\Pi$. If
\Cref{step:pd-comp-ph} used a fixed threshold value, then there would exist
pseudo-deterministic $\mif$ algorithms yielding element frequencies close to
this threshold, where even low-magnitude noise could affect which elements are
included in $P_h$.  Using $C$ to choose random threshold values that are
independent of the choice of $A$ ensures that most of the time, the noise has
no influence on $P_h$; ultimately, ensuring that $\fco(\Pi,C,x,k)$ and
$\fco(A,C,x,k)$ most likely produce the same output. More detail is given in
the proof of \Cref{lem:pd-stability}.

Second, the design of \Cref{alg:pd-set-finding} ensures that
$\fco(\Pi,C,\emptystream,d)$
actually produces a set of $w_d$ possible outputs. Recall the $\avoid(m,a,b)$
communication problem described in \Cref{subsec:rt-lb-overview}. The output
sets $T_y$ of size $w_{d-1}$ computed on \Cref{step:pd-compute-t}, for
each $y \in \SSD{S}{t_d}$, will have a similar structure to the input
and output sets for an $\avoid$ protocol, in that the set $T_y$ is typically
disjoint from $y$, and in that the distribution of possible $T_y$ values is
limited (through the requirement that $\Pi$ agrees with a mixture of functions
corresponding to deterministic streaming algorithms). Note that we
cannot just extract the ``most common values'' that occur in the sets $T_y$
for $y \in \SSD{S}{t_d}$, because if $t_d w_{d-1} \ll |S|$, it is
possible that these are very concentrated; for example there could exist a set
$V$ of size $w_{d-1}$ so that if $y \cap V = \emptyset$, then $T_y = V$, in
which case an algorithm satisfying the first property can only reliably
identify $w_{d-1}$ outputs. Instead, we iteratively build up a set of common
outputs, using the following observation. If the current set of common
algorithm outputs $Q_{h-1}$ is smaller than $w_d$, then the set $P_h$
containing the most common elements of $T_y$ for $y \in
\SSD{Q_{h-1}}{t_d}$ can not be entirely contained by $Q_{h-1}$: if we
did have $P_h \subseteq Q_{h-1}$, then one can show it is possible to
construct an impossibly efficient protocol for
$\avoid(|Q_{h-1}|,t_d,w_{d-1})$. Consequently, until for some value of $h$ we
have $|Q_{h}| \ge w_d$, it is possible to find a slightly larger $Q_{h+1}$.
More detail is given in the proof of \Cref{lem:pd-set-growth}.

\medskip Before beginning the formal proof, let us set the various parameters
precisely. Concretely, we set
\begin{align}
  p &= \ceil{\max\left(\sqrt{\frac{10 \rl \log (64|S|)}{3 z \log \frac{1}{2\delta}}}, 
    \frac{30 \log(64|S|)}{\log \frac{1}{2\delta}}\right) } \,.  \label{eq:pd-p-choice} \\
  d &= 1 + \floor{\frac{\rl}{18 zp}} \,, \label{eq:d-def} \\
  t_d = t_{d-1} = \cdots = t_2 &= \ceil{4 \ln 2 (zp+2)} \,. \\
\intertext{Recall also that}
  \epsilon &\le (2\delta)^{p/30} \,. \\
\intertext{We further define}
  \epsilon_k &:= w_k (64 |S|)^{k-1} \epsilon \,, \label{eq:eps-k-def}
\end{align}
which will upper-bound the probability that $\fco(A, C, x, k) \ne \fco(\Pi, C, x, k)$.
%
We record some useful estimations in the next lemma.

\begin{lemma}\label{lem:pd-p-properties}
  With the above choices, $t_d = \cdots = t_2 \le 9zp$,\, $t_1 \ge \rl / 2$,\, and
  $\epsilon \le 1 / (64 |S|)^d$.
\end{lemma}
%
\begin{proof}%[Proof of \Cref{lem:pd-p-properties}]
  The claim about the $t_i$s holds since $z \ge 1$ and $p \ge 1$. Using this
  in \cref{eq:d-def} gives $t_1 \ge \rl - (d - 1) 9zp \ge \rl / 2$. We turn to
  proving the claim about $\epsilon$. Recall that, by design, $\epsilon \le
  (2\delta)^{p/30}$. 

  The two branches of the maximum in \cref{eq:pd-p-choice} ensure that:
  \begin{align}
    p^2 \ge \frac{10 \rl \log (64|S|)}{3 z \log \frac{1}{2\delta}} \qquad \text{and} \qquad p \ge \frac{30 \log(64|S|)}{\log \frac{1}{2\delta}}\,. \label{eq:pd-def-p}
  \end{align}
  Because $d \le \max(1, \frac{\rl}{9 z p})$, it suffices to prove that
  $\epsilon \le 1 / (64 |S|)$ and that $\epsilon \le 1 / (64 |S|)^{\rl / (9 z
  p)}$. Now,
  \begin{align*}
    \log \frac{1}{\epsilon}
    &\underset{\text{\Cref{lem:error-reduction-by-vote}}}{\ge} \frac{p}{30}
      \log \frac{1}{2\delta} \underset{\text{\cref{eq:pd-def-p}}}{\ge}
      \frac{1}{30 p} \frac{10 \rl \log(64 |S|)}{3z} = \frac{\rl}{9zp}
      \log(64|S|) \,; \\ 
    \log \frac{1}{\epsilon} 
    &\underset{\text{\Cref{lem:error-reduction-by-vote}}}{\ge} \frac{p}{30} 
      \log \frac{1}{2\delta} \underset{\text{\cref{eq:pd-def-p}}}{\ge} \log(64|S|) \,. \qedhere
  \end{align*}
\end{proof}


\subsection{Common Outputs Behave Canonically}

We now come to the central lemma in the proof, which asserts that the set of
common outputs is likely the same for the canonical function $\Pi$ as it is
for a random draw $A \sim \cB$. It also asserts two other key properties of
\fco. The lemma can be thought of as a ``proof of correctness'' of \fco.

\begin{lemma}\label{lem:pd-main-recursion}
  Let $k \in [d]$ and $x \in [n]^{t_d + \cdots + t_{k+1}}$. Then \fco
  satisfies the following properties.
  \begin{enumerate}
    \item $\Pr_{A \sim \cB, C \in_R [1,2)^{d \times \NN}}[\fco(A, C, x, k) =
    \fco(\Pi, C, x, k)] \ge 1 - \epsilon_k$.
    \item For all $C \in [1,2)^d$, the set $\fco(\Pi, C, x, k)$ is disjoint from $x$ and a subset of $S$. 
    \item For all $A \colon [n]^\rl \to [n]$ and $C \in [1,2)^d$,
    $\fco(A,C,x,k)$ outputs a set of size $w_k$.
  \end{enumerate}
\end{lemma}
\begin{proof}
  The proof is by induction on $k$, spread over the next few lemmas. The case
  $k = 1$ is handled in \Cref{lem:pd-base-case}. For the induction step, the
  heart of the argument, which invokes a communication lower bound for \avoid, is
  given in \Cref{lem:pd-set-growth}. Following this,
  \Cref{lem:pd-step-case-aux} establishes the latter two claims in the lemma
  and \Cref{lem:pd-stability} establishes the first claim.
\end{proof}

\begin{lemma}\label{lem:pd-base-case} 
  \Cref{lem:pd-main-recursion} holds for $k = 1$.
\end{lemma}
%
\begin{proof}%[Proof of \Cref{lem:pd-base-case}]
  Let $e_0,\ldots,e_{t_1}$ be the values of the variables on Lines
  \ref{step:pd-base-case-start} to \ref{step:pd-base-case-mid} of
  \Cref{alg:pd-set-finding} when $\fco(\Pi, C, x, 1)$ is called; note that
  these do not depend on $C$. For $i \in \{0,\ldots,t_1\}$, define sequences
  $s_i = \langle e_0,\ldots,e_{i-1},1,\ldots,1 \rangle$, so that $s_0 =
  \langle 1,1,\ldots,1 \rangle$, and $s_{t_1} = \langle e_0,\ldots,e_{t_1-1}
  \rangle$. If, for all $i \in \{0,\ldots,t_1\}$, we have $A(x \concat s_i) =
  \Pi(x \concat s_i)$, the value of $\fco(A, C, x, 1)$ will exactly match
  $\fco(\Pi,C,x,1)$. By a union bound,
  \begin{align*}
    \Pr_{A \sim \cB, C}[\fco(A, C, x, k) \ne \fco(\Pi, C, x, k)] 
    \le \sum_{i = 0}^{t_1} \Pr_{A \sim \cB}[A(x \concat s_i) \ne \Pi(x \concat s_i)] 
    \underset{\text{\cref{eq:pd-fn-prop}}}{\le} (t_1 + 1) \epsilon 
    = \epsilon_1 \,.
  \end{align*}
  
  Because $\Pi$ is the canonical output function for a protocol for $\mif$,
  for any $z \in [n]^\rl$, we have $\Pi(z) \notin z$. Consequently, each $e_i
  = \Pi(x \concat \langle e_0, \ldots, e_{i-1}, 1, \ldots, 1\rangle)$ is
  neither contained in $x$ nor by $\{e_0,\ldots,e_{i-1}\}$; thus
  $\{e_0,\ldots,e_{t_1}\}$ has size $t_1 + 1 = w_1$ and is disjoint from $x$.
\end{proof}

\begin{lemma}\label{lem:pd-set-growth}
  Let $x \in [n]^{t_d+\cdots+t_{k+1}}$. 
  When computing $\fco(\Pi,C,x,k)$, in the $h$th loop iteration, if $|Q_{h -
  1}| < w_k$, then $|P_h \setminus Q_{h-1}| \ge \ceil{\frac{1}{4} w_{k-1}}$.
  Consequently, the algorithm will return using \Cref{step:pd-good-exit},
  not \Cref{step:pd-bad-exit}.
\end{lemma}
%
\begin{proof}%[Proof of \Cref{lem:pd-set-growth}]
  Assume for sake of contradiction that $|Q_{h - 1}| < w_k$ and $|P_h
  \setminus Q_{h-1}| \le \floor{\frac{1}{4} w_{k-1}}$. Then we can use the
  algorithm $\cA$ to implement a protocol for the one-way communication
  problem $\avoid(|Q_{h-1}|, t_k, \ceil{\frac{1}{2} w_{k-1}})$, with $\le
  \frac{1}{2}$ probability of error.
  
  We assume without loss of generality that $Q_{h-1} = [|Q_{h-1}|]$; if not,
  relabel coordinates so that this holds. In the protocol, after Alice is
  given a subset $W \subseteq Q_{h-1}$ with $|W| = t_k$, they construct a
  sequence $v = x \concat \sort(W)$ in $[n]^{t_d + \cdots + t_k}$.  Then Alice
  uses public randomness to instantiate an instance $E$ of $\cA$;  inputs the
  sequence $v$ to $E$; and sends the new state of $E$ to Bob, using a $z
  p$-bit message. As Bob shares the public randomness,  they can use this
  state to evaluate the output of the algorithm on any continuation of the
  stream. In particular, Bob can evaluate the algorithm for any possible
  suffix, to produce a function $\tilde{A}_{x \concat \sort(W)} :
  [n]^{t_{k-1}+\cdots+t_1} \rightarrow [n]$; Bob then samples a random $C \in
  [1,2)^{d \times \NN}$, and computes $V = \fco(\tilde{A}_{x \concat
  \sort(W)}, C, k-1)$, which is a subset of $S$. If $|V \cap Q_{h-1}| \ge
  \ceil{\frac{1}{2} w_{k-1}}$, Bob outputs the smallest $\ceil{\frac{1}{2}
  w_{k-1}}$ entries of $V \cap Q_{h-1}$. Otherwise, Bob outputs an arbitrary
  set of size $\ceil{\frac{1}{2} w_{k-1}}$.
  
  First, we observe that for any value of $\sort(W)$, the distribution of
  $\tilde{A}_{x \concat \sort(W)}$
  is exactly the same as the distribution of $A_{x \concat \sort(W)}$, when $A$ is drawn from $\cB$;
  this follows because for a fixed setting of the oracle random string of the algorithm,
  it behaves deterministically.
  
  Applying \Cref{lem:pd-main-recursion} at $k-1$, we observe that for any $W \in \binom{Q_{h-1}}{t_k}$,
  \begin{align*}
    \Pr[\fco(\tilde{A}_{x \concat \sort(W)}, C, k-1) = \fco(\Pi, C, x \concat \sort(W), k-1)] 
    \ge 1 - \epsilon_{k-1} \ge \frac{3}{4} \,.
  \end{align*}
  Furthermore, we are guaranteed that $\fco(\Pi, C, x \concat \sort(W), k-1)$
  has size $w_{k-1}$ and is disjoint from $W$.
  
  We now bound the probability, over a uniformly random $y \in
  \SSD{Q_{h-1}}{t_k}$, that $|\fco(\Pi, C, x \concat y, k-1) \cap
  Q_{h-1}| < \ceil{\frac{1}{2} w_{k-1}}$. Define $T_y = \fco(\Pi, C, x \concat y, k-1)$ and, for each $j\in S$,
  $f_j^{(h)}$, as in \Cref{alg:pd-set-finding}. In particular, we have:
  \begin{align}
    \Pr_{y,C}\left[|T_y \cap Q_{h-1}| < \ceil{\frac{1}{2} w_{k-1}} \right]
      &= \Pr_{y,C}\left[|T_y \setminus Q_{h-1}| > \floor{\frac{1}{2} w_{k-1}} \right] \nonumber \\
      &\le \Pr_{y,C}\left[|T_y \setminus P_h \setminus Q_{h-1}| > \floor{\frac{1}{2} w_{k-1}} - \floor{\frac{1}{4} w_{k-1}} \right] \label{eq:pd-lb-minor-set-drop} \\
      &\le \Pr_{y,C}\left[|T_y \setminus P_h| \ge \frac{1}{4} w_{k-1} \right] \,. \nonumber
  \end{align}
  (The inequality on \cref{eq:pd-lb-minor-set-drop} follows since we
  assumed $|P_h \setminus Q_{h-1}| \le \floor{\frac{1}{4} w_{k-1}}$.) Note that:
  \begin{align}
    \sum_{j \notin P_h} f_j^{(h)} = \sum_{y \in \SSD{Q_{h-1}}{t_k}} |T_y \setminus P_h| \ge \frac{1}{4} w_{k-1} \left|\left\{y \in \SSD{Q_{h-1}}{t_k} : |T_y \setminus P_h| \ge \frac{1}{4} w_{k-1}\right\}\right| \label{eq:pd-lb-count-part} \,.
  \end{align}
  Using the fact that $y$ is uniformly distributed over $\SSD{Q_{h-1}}{t_k}$, gives:
  \begin{align*}
    \Pr_{y,C}\left[|T_y \setminus P_h| \ge \frac{1}{4} w_{k-1} \right]
      &= \EE_C \frac{\left|\left\{y \in \SSD{Q_{h-1}}{t_k} : |T_y \setminus P_h| \ge \frac{1}{4} w_{k-1}\right\}\right|}{\binom{|Q_{h-1}|}{t_k}} \\ 
      &\le \EE_C \frac{\sum_{j \notin P_h} f_j^{(h)}}{\frac{1}{4} w_{k-1} \binom{|Q_{h-1}|}{t_k}} && \text{by \cref{eq:pd-lb-count-part}}\\ 
      &\le \EE_C \frac{(|S|-|P_h|) \frac{C_{k,h} w_{k-1}}{16 |S|} \binom{|Q_{h-1}|}{t_k} }{\frac{1}{4} w_{k-1} \binom{|Q_{h-1}|}{t_k}} &&\text{by definition of $P_h$} \\
      &= \EE_C \frac{C_{k,h}}{4} \frac{|S|-|P_h|}{|S|} \le \frac{1}{2} \,. &&\text{since $|P_h|\ge 0$, $C_{k,h} \le 2$} 
  \end{align*}
  Thus the probability that $|T_y \cap Q_{h-1}| < \ceil{\frac{1}{2} w_{k-1}}$
  holds is $\le 1/2$. Since Bob only gives an incorrect output when this
  happens or when $\fco(\tilde{A}_{x \concat \sort(W)}, C, k-1) \ne
  \fco(\Pi, C, x \concat \sort(W), k-1)$, it follows by a union bound that
  the total failure probability is $\le \frac{1}{2} + \frac{1}{4} \le
  \frac{3}{4}$.
  
  Consequently, the protocol implementation has $\le \frac{3}{4}$ error when inputs are drawn from the uniform distribution over $\binom{Q_{h-1}}{t_k}$; by \Cref{lem:avoid-lb}, we 
  obtain a lower bound on the required message length, giving
  \begin{align*}
    z p > \frac{t_k \ceil{\frac{1}{2} w_{k-1}}}{|Q_{h-1}| \ln 2} + \log(1 - 3/4) \ge \frac{t_k w_{k-1}}{|Q_{h-1}| \cdot 2 \ln  2} - 2 \,.
  \end{align*}
  Rearranging this slightly and using integrality of $|Q_{h-1}|$ gives:
  \begin{align*}
    |Q_{h-1}| \ge \ceil{\frac{t_k w_{k-1}}{2 \ln  2 (zp+2) }} = \ceil{\frac{\ceil{4 \ln 2 (zp + 2)}}{2 \ln  2 (zp+2) } w_{k-1}} \ge 2 w_{k-1} = w_k \,,
  \end{align*}
  but as $|Q_{h-1}| < w_k$, this implies $w_k < w_k$, which is a
  contradiction; this proves that the assumption $|P_h \setminus Q_{h-1}| \le
  \frac{1}{4} w_{k-1}$ must have been invalid.
  
  Finally, we observe that since, in each iteration of the loop on Lines \ref{step:pd-for-loop-start} to \ref{step:pd-for-loop-end},
  $|Q_h| = |Q_{h-1} \cup P_{h}| = |Q_{h-1}| + |P_{h} \setminus Q_{h-1}| \ge |Q_{h-1}| + \ceil{\frac{1}{4} w_{k-1}}$,
  and we initially have $|Q_0| = w_{k-1}$, the size of $Q_h$ (assuming we haven't returned yet)
  must be $\ge w_{k-1} (1 + h / 4)$. By the last loop iteration (with $h=4$), we will have $|Q_h| \ge 2 w_{k-1} = w_k$.
\end{proof}


\begin{lemma}\label{lem:pd-step-case-aux}
  For $k > 1$, $x \in [n]^{t_d+\cdots+t_{k+1}}$, $\fco(\Pi, C, x, k)$ is
  disjoint from $x$ and a subset of $S$; and for all $A, C, k$,
  $\fco(A,C,x,k)$ outputs a set of size $w_k$.
\end{lemma}
%
\begin{proof}%[Proof of \Cref{lem:pd-step-case-aux}]
  By \Cref{lem:pd-main-recursion} at $k - 1$, the sets $T_{A,x \concat y}$
  chosen on \Cref{step:pd-compute-t} are always subsets of $S$ and
  disjoint from $x \concat y$, and hence disjoint from $x$. Per
  \Cref{lem:pd-set-growth}, \textsc{FindCommonOutputs} will return a subset of
  $Q_h$ using \Cref{step:pd-good-exit}, where $h$ is the last loop
  iteration number. Each element of $Q_h$ was either in $T_{A,x \circ \langle
  1,2,\ldots,t_k\rangle}$ (and hence also in $S$) or was in $P_{h'}$ for some
  $h' \le h$. Note that $P_{h'}$ only contains integers $j$ for which
  $f_j^{(h')} > 0$; i.e., which were contained in one of the sets $(T_{A,x
  \concat y})_{y \in \SSD{Q_{h'-1}}{t_k}}$, and are thereby also in $S$.
  Consequently, the set returned is contained in $S$, which implies $|S| \ge
  w_k$.

  Calls to $\fco(A,C,x,k)$ will either output through
  \Cref{step:pd-good-exit} (where the size of the set has been checked by
  the pseudocode) or through \Cref{step:pd-bad-exit} (where a subset of
  size $w_k$ must exist, because we know $|S| \ge w_k$).
\end{proof}

\begin{lemma}\label{lem:pd-stability}
  For $k > 1$, and all $x \in [n]^{t_d + \cdots + t_{k+1}}$, 
  \begin{align*}
    \Pr_{A \sim \cD, C}[\fco(A, C, x, k) \ne \fco(\Pi, C, x, k)] \le \epsilon_k \,.
  \end{align*}
\end{lemma}
%
\begin{proof}%[Proof of \Cref{lem:pd-stability}]
  The proof of the lemma follows from the observation that, when computing
  $\fco(A,C,x,k)$, even if a fraction of the recursive calls to $\fco(A,C,x
  \concat y,k-1)$ produced incorrect outputs, the values for $Q_0$ and $(P_h)_{h \ge 1}$ will likely match those computed when $\fco(\Pi,C,x,k)$ is called.
  
  Henceforth, we indicate variables from the computation of $\fco(\Pi,C,x,k)$ without a tilde,
  and variables from the computation of $\fco(A,C,x,k)$ with a tilde. For example, $f_j^{(h)}$
  is computed using $B = \Pi$, while $\tilde{f}_j^{(h)}$ is computed using $B = A$. We also define
  \begin{align*}
    \hf_j^{(h)} &= \left|\left\{ y \in \SSD{Q_{h-1}}{t_k} : j \in T_{A,y} \right\}\right| \\
    \hP_h &= \left\{j \in S : \hf_j^{(h)} \ge \frac{C_{k,h} w_{k-1}}{16 |S|} \left|\binom{Q_{h-1}}{t_k}\right| \right\} \,;
  \end{align*}
  that is, $\hf_j^{(h)}$ and $\hP_h$ are the values that would be computed by $\fco(A,C,x,k)$ if the set $Q_{h-1}$ was used instead of the set $\tQ_{h-1}$.
  
  Say $\fco(\Pi, C, x, k)$ returns from the loop at iteration $h^\star$. The output of
  $\fco(A, C, x, k)$ will equal $\fco(\Pi, C, x, k)$ if $Q_0 = \tQ_0$ and for all $h \in [h^\star]$,
  we have $P_h = \hP_h$. (If this occurs, then as $Q_0 = \tQ_0$, $\hP_1 = \tP_1$, so
  $Q_1 = Q_0 \cup P_1 = \tQ_0 \cup \tP_1 = \tQ_1$, and as $Q_1 = \tQ_1$, $\hP_2 = \tP_2$, and so on.)  By \Cref{lem:pd-main-recursion} at $k-1$, the probability that
  $Q_0 \ne \tQ_0$ is $\le \epsilon_{k-1}$. Consider a specific $h \in [h^\star]$;
  the only way in which $\hP_h \ne P_h$ is if there is some $j \in S$ for which $f_j^{(h)}$ and $\hf_j^{(h)}$
  are on opposite sides of the threshold $\frac{C_{k,h} w_{k-1}}{16 |S|} |\binom{Q_{h-1}}{t_k}|$.
  
  Let $\lambda_{h}$ be the random variable indicating the fraction of $y \in \SSD{Q_{h-1}}{t_k}$
  for which $T_{A,x \concat y} \ne T_{\Pi,x \concat y}$. Note that the values $T_{A,x \concat y}$ are functions of
  the random variable $A$ and of $C_{k',h}$ for $k' < k, h \in \NN$; in particular $T_{A,x\circ y}$ is independent of $(C_{k,h})_{h \in \NN}$. By \Cref{lem:pd-main-recursion} at $k-1$,
  $\Pr[T_{A,x\circ y} \ne T_{\Pi,x\circ y}] \le \epsilon_{k-1}$, which implies $\EE \lambda_h \le \epsilon_{k-1}$.
  
  Fix a particular setting of $A$ and $(C_{k',h})_{k' < k, h \in \NN}$. Since each set $T_{A,x\circ y}$
  contributes $1$ unit to each of $w_{k-1}$ variables $\hf_j^{(h)}$:
  \begin{align*}
    \sum_{j \in S} \left|f_j^{(h)} - \hf_j^{(h)}\right| \le w_{k-1} \left|\left\{y \in \SSD{Q_{h-1}}{t_k} : T_{A,x \concat y} \ne T_{\Pi,x\circ y}\right\}\right| = w_{k-1} \lambda_h \binom{Q_{h-1}}{t_k} \,.
  \end{align*}
  Let $F$ be the set of possible values in $[1,2)$ for $C_{k,h}$ for which $P_h \ne \hP_h$;
  this is a union of intervals corresponding to each pair $\left(f_j^{(h)}, \hf_j^{(h)}\right)$,
  for $j \in S$. A given value $c$ is bad for $j$ if 
  \begin{align*}
    f_j^{(h)} < \frac{c w_{k-1}}{16 |S|} \binom{|Q_{h-1}|}{t_k} \le \hf_j^{(h)} \,; \qquad \text{equivalently:} \qquad c \in \left(\frac{16 |S| f_j^{(h)} }{w_{k-1} \binom{|Q_{h-1}|}{t_k}}, \frac{16 |S| \hf_j^{(h)}}{w_{k-1} \binom{|Q_{h-1}|}{t_k}}\right] \,,
  \end{align*}
  and similarly in the case where $\hf_j^{(h)} < f_j^{(h)}$. The measure of $F$ is:
  \begin{align*}
    \le \sum_{j \in S} \frac{16 |S|}{w_{k-1} \binom{|Q_{h-1}|}{t_k}} |\hf_j^{(h)} - f_j^{(h)}| \le \frac{16 |S|}{w_{k-1}} w_{k-1} \lambda_h = 16 |S| \lambda_h \,.
  \end{align*}
  This upper bounds the probability that $C_{k,h} \in F$ and $P_h \ne \hP_h$. We then have:
  \begin{align*}
    \Pr[P_h \ne \hP_h] &= \EE_{A, (C_{k',h})_{k' < k}} \Pr[ C_{k,h} \in F ] 
      \le \EE_{A, (C_{k',h})_{k' < k}} (16 |S| \lambda_h) = 16 |S| \epsilon_{k-1} \,.
  \end{align*}
  
  By a union bound, the probability that $Q_0 \ne \tQ_0$ or $P_h \ne \hP_h$ for any $h \le h^\star$ is
  \begin{align*}
    \le \epsilon_{k-1} + h^\star 16 |S| \epsilon_{k-1} \le (1 + 4 \cdot 16 |S|) \epsilon_{k-1} \le \frac{64 |S| w_k}{w_{k-1}} \epsilon_{k-1} \,.
  \end{align*}
  Thus $\Pr[\fco(A, C, x, k) \ne \fco(\Pi, C, x, k)] \le \frac{64 |S| w_k}{w_{k-1}} \epsilon_{k-1} = \epsilon_k$.
\end{proof}

We have thus completed the proof of \Cref{lem:pd-main-recursion}, establishing
the key properties of \fco. It is time to use them to derive our lower bound.

\subsection{Obtaining a Pseudo-Deterministic Lower Bound}

\begin{lemma}\label{lem:pd-z-lb} 
  We have:
  \begin{align*}
    z \ge \frac{\rl}{\log\frac{2|S|}{\rl}} \min\left(\frac{1}{36},
    \frac{\log\frac{1}{2\delta}}{17280 \log(64|S|) \log\frac{2|S|}{\rl}} \right) \,.
  \end{align*}
\end{lemma}
%
\begin{proof}%[Proof of \Cref{lem:pd-z-lb}]
  A consequence of \Cref{lem:pd-main-recursion} is that $\fco(\Pi, C, d)$ will
  output a set of size $w_d$ which is a subset of $S$. This shows that $|S|
  \ge w_d$.  Now, from the definition of $w_d$, it follows that
  \begin{align*}
    |S| \ge w_d = 2^{d-1} (t_1 + 1) > 2^{d-1} t_1 \ge 2^{d-1} \frac{\rl}{2} \qquad \implies \qquad \log \frac{2 |S|}{\rl} \ge d - 1 = \floor{\frac{\rl}{18 z p}} \,.
  \end{align*}
  Since $|S| \ge \rl + 1$, the left hand side $\log \frac{2 |S|}{\rl} > 1$, so using the inequality $x/2 \le \max(\floor{x}, 1)$ gives:
  \begin{align*}
    \frac{\rl}{36 z p} \le \log \frac{2 |S|}{\rl} \qquad \implies \qquad z \ge \frac{\rl}{36 \log \frac{2|S|}{\rl}} \cdot \frac{1}{p} \,.
  \end{align*}
  Next, we expand the definition of $p$ (see \cref{eq:pd-p-choice}), eliminating the ceiling using the inequality $\ceil{x} \le \max(1, 2 x)$:
  \begin{align*}
    z &\ge \frac{\rl}{36 \log \frac{2|S|}{\rl}} \min\left(1, \frac{1}{2} \sqrt{ \frac{3 z \log \frac{1}{2\delta}}{10 \rl \log(64 |S|) } }, \frac{1}{2} \cdot \frac{\log \frac{1}{2\delta}}{30 \log(64 |S|) }\right) \,.
  \end{align*}
  We have two cases: if the left or right side of the minimum is smallest, then:
  \begin{align}
    z \ge \min\left( \frac{\rl}{36 \log \frac{2|S|}{\rl}}, \frac{\rl \log \frac{1}{2\delta}}{2160 \log \frac{2|S|}{\rl} \log(64|S|)} \right) \,, \label{eq:pd-min-lr}
  \end{align}
  while otherwise, if the center is smallest, we get:
  \begin{align*}
    z^2 \ge  \frac{1}{4} \frac{\rl^2}{(36 \log \frac{2|S|}{\rl})^2} \cdot \frac{3 z \log \frac{1}{2\delta}}{10 \rl \log(64 |S|) } \qquad \implies \qquad z \ge \frac{\rl \log \frac{1}{2 \delta}}{17280 (\log \frac{2|S|}{\rl})^2 \log(64|S|)} \,.
  \end{align*}
  As $\log\frac{2|S|}{\rl} \ge 1$, this is smaller than the right minimum
  branch of \cref{eq:pd-min-lr}, so the common lower bound for all cases is:
  \begin{align*}
    z &\ge \min\left( \frac{\rl}{36 \log \frac{2|S|}{\rl}}, \frac{\rl \log \frac{1}{2 \delta}}{17280 (\log \frac{2|S|}{\rl})^2 \log(64|S|)} \right) \,. \qedhere
  \end{align*}
\end{proof}
% 


\begin{lemma}\label{lem:pd-s-le-2zp1}
  We have $|S| < 2^{z+1}$.
\end{lemma}
%
\begin{proof}%[Proof of \Cref{lem:pd-s-le-2zp1}]
  For each $a \in S$, let $x_a \in \Pi^{-1}(a)$. One can use $\cA$ to provide a randomized $\le\delta$-error, $z$-bit encoding of the elements in $S$. Using public randomness, encoder and decoder choose the oracle random string for $\cA$. Each $a \in S$ is encoded by sending $x_a$ to $\cA$ and outputting the algorithm state $\sigma$. To decode, given a state $\sigma$, one evaluates the output of $\cA$ at state $\sigma$. Using the minimax principle, one can prove that the randomized encoding requires $\ge \log( (1-\delta) |S| )$ bits of space, which implies $2^z \ge (1-\delta) |S|$. Since $\delta \le \frac{1}{3}$, it follows $s \le \frac{3}{2} 2^{z} < 2^{z+1}$.
\end{proof}

We now establish the main result.


\begin{theorem} %{pdlbintro}
\label{thm:pd-lb-intro}
  Pseudo-deterministic $\delta$-error random oracle algorithms for $\mif(n,\rl)$ require
  \begin{align*}
    \Omega\left(\min\left(\frac{\rl}{\log \frac{2n}{\rl}} + \sqrt{\rl}, \frac{\rl \log \frac{1}{2\delta}}{(\log \frac{2n}{\rl})^2 \log n} + \left(\rl \log \frac{1}{2\delta}\right)^{1/4}\right)\right)
  \end{align*}
  bits of space when $\delta \le \frac{1}{3}$. In particular, when $\delta = 1/\poly(n)$ and $\rl = \Omega(\log n)$, this is:
  \begin{align*}
    \Omega\left(\frac{\rl}{(\log\frac{2 n}{\rl})^2} + \left(\rl \log n\right)^{1/4}\right) \,.
  \end{align*}
\end{theorem}

\begin{proof}[Proof sketch]
  Using \Cref{lem:pd-s-le-2zp1} and the fact that $S \subseteq [n]$, we obtain
  $|S| \le \min(n, 2^{z+1})$. The theorem follows by combining this bound with
  the inequality of \Cref{lem:pd-z-lb}, and for each of four cases
  corresponding to different branches of $\min$ and $\max$, solving to find a
  lower bound on $z$. The full proof with calculations is given in
  \Cref{subsec:pd-lb-mech}.
\end{proof}

\subsection{Mechanical Proofs for \texorpdfstring{\Cref{sec:pd-lb}}{Pseudo-Deterministic and Random Seed Lower Bounds}}\label{subsec:pd-lb-mech}


\begin{proof}[Proof of \Cref{thm:pd-lb-intro}]
  By \Cref{lem:pd-z-lb}, we have:
  \begin{align}
    z \ge \min\left(
      \frac{\rl}{36 \log\frac{2 |S|}{\rl}},
      \frac{\rl}{17280 \log\frac{2 |S|}{\rl}} \cdot \frac{\log(1/2\delta)}{\log(64 |S|) \log\frac{2 |S|}{\rl}}
    \right) \,.\label{eq:pd-lb-restate}
  \end{align}
  By \Cref{lem:pd-s-le-2zp1}, $|S| \le \min(n, 2^{z+1}) \le \min(n, 4^z)$. We will apply this
  inequality to each branch of the minimum in Eq. \ref{eq:pd-lb-restate}. First, say that
  the left part of the minimum is larger than the right. Then $z \ge \rl / (36 \log\frac{2 |S|}{\rl})$.
  Applying $|S| \le n$ and $|S| \le 4^z$, this implies:
  \begin{align*}
    z &\ge \frac{\rl}{36 \log\frac{2 n}{\rl}} \,, \qquad \text{and} \\
    z &\ge \frac{\rl}{36 \log\frac{2 \cdot 4^z}{\rl}} \ge \frac{\rl}{36 \cdot 3 z} \qquad \implies \qquad z \ge \sqrt{\frac{\rl}{108}} \,.
  \end{align*}
  Thus:
  \begin{align}
    z \ge \max\left(\frac{\rl}{36 \log\frac{2 n}{\rl}}, \sqrt{\frac{\rl}{108}} \right) \,.\label{eq:pd-lb-case-l}
  \end{align}

  Next, say that the right side of the minimum in Eq. \ref{eq:pd-lb-restate} is larger.
  Then applying $|S| \le n$ and $|S| \le 4^z$ to that side, we get:
  \begin{align*}
    z &\ge \frac{\rl \log(1/2\delta)}{17280 (\log\frac{2 n}{\rl})^2 \log(64 n)} \,, \qquad \text{and} \\
    z &\ge \frac{\rl \log(1/2\delta)}{17280 (\log\frac{2 \cdot 4^z}{\rl})^2 \log(64 \cdot 4^z)} \ge \frac{\rl \log(1/2\delta)}{17280 \cdot 3^2 \cdot 8 z^3} \qquad \implies \qquad z \ge \left(\frac{\rl \log(1/2\delta)}{1244160}\right)^{1/4} \,.
  \end{align*}
  Thus:
  \begin{align}
    z \ge \max\left(
      \frac{\rl \log(1/2\delta)}{17280 (\log\frac{2 n}{\rl})^2 \log(64 n)},
      \left(\frac{\rl \log(1/2\delta)}{1244160}\right)^{1/4}
    \right) \,.\label{eq:pd-lb-case-r}
  \end{align}

  The minimum of the lower bounds from Eqs. \ref{eq:pd-lb-case-l} and \ref{eq:pd-lb-case-r}
  holds in all cases, so:
  \begin{align*}
    z \ge \min\left(\max\left(\frac{\rl}{36 \log\frac{2 n}{\rl}}, \sqrt{\frac{\rl}{108}}\right), \max\left(\frac{\rl \log(1/2\delta)}{17280 (\log\frac{2 n}{\rl})^2 \log(64n)},\left(\frac{\rl \log(1/2\delta)}{1244160}\right)^{1/4}\right)\right) \,.
  \end{align*}
\end{proof}

\newcommand{\etalchar}[1]{$^{#1}$}
\begin{thebibliography}{GGMW20}

\bibitem[ABJ{\etalchar{+}}22]{AjtaiBJSSWZ22}
Mikl\'os Ajtai, Vladimir Braverman, T.S. Jayram, Sandeep Silwal, Alec Sun,
  David~P. Woodruff, and Samson Zhou.
\newblock The white-box adversarial data stream model.
\newblock In {\em Proc. 41st ACM Symposium on Principles of Database Systems},
  pages 15--27, 2022.

\bibitem[ACGS23]{AssadiCGS23}
Sepehr Assadi, Amit Chakrabarti, Prantar Ghosh, and Manuel Stoeckl.
\newblock Coloring in graph streams via deterministic and adversarially robust
  algorithms.
\newblock In {\em Proc. 42nd ACM Symposium on Principles of Database Systems},
  pages 141--153, 2023.

\bibitem[BEEO22]{BenEliezerEO22}
Omri Ben-Eliezer, Talya Eden, and Krzysztof Onak.
\newblock Adversarially robust streaming via dense-sparse trade-offs.
\newblock In {\em Symposium on Simplicity in Algorithms (SOSA)}, pages
  214--227, 2022.

\bibitem[BJWY20]{BenEliezerJWY20}
Omri Ben{-}Eliezer, Rajesh Jayaram, David~P. Woodruff, and Eylon Yogev.
\newblock A framework for adversarially robust streaming algorithms.
\newblock In {\em Proc. 39th ACM Symposium on Principles of Database Systems},
  page 63–80, 2020.

\bibitem[BKKS23]{BravermanKKS23}
Vladimir Braverman, Robert Krauthgamer, Aditya Krishnan, and Shay Sapir.
\newblock Lower bounds for pseudo-deterministic counting in a stream.
\newblock {\em arXiv preprint arXiv:2303.16287}, 2023.

\bibitem[BR93]{BellareR93}
Mihir Bellare and Phillip Rogaway.
\newblock Random oracles are practical: A paradigm for designing efficient
  protocols.
\newblock In {\em Proc. 1st ACM Conference on Computer and Communications
  Security}, pages 62--73, 1993.

\bibitem[BY20]{BenEliezerY20}
Omri Ben{-}Eliezer and Eylon Yogev.
\newblock The adversarial robustness of sampling.
\newblock In {\em Proc. 39th ACM Symposium on Principles of Database Systems},
  pages 49--62. {ACM}, 2020.

\bibitem[CGS22]{ChakrabartiGS22}
Amit Chakrabarti, Prantar Ghosh, and Manuel Stoeckl.
\newblock Adversarially robust coloring for graph streams.
\newblock In {\em Proc. 13th Conference on Innovations in Theoretical Computer
  Science}, pages 37:1--37:23, 2022.

\bibitem[Fei19]{Feige19}
Uriel Feige.
\newblock A randomized strategy in the mirror game.
\newblock {\em arXiv preprint arXiv:1901.07809}, 2019.

\bibitem[GGMW20]{GoldwasserGMW20}
Shafi Goldwasser, Ofer Grossman, Sidhanth Mohanty, and David~P. Woodruff.
\newblock {Pseudo-Deterministic Streaming}.
\newblock In {\em Proc. 20th Conference on Innovations in Theoretical Computer
  Science}, volume 151, pages 79:1--79:25, 2020.

\bibitem[GGS23]{GrossmanGS23}
Ofer Grossman, Meghal Gupta, and Mark Sellke.
\newblock Tight space lower bound for pseudo-deterministic approximate
  counting.
\newblock {\em arXiv preprint arXiv:2304.01438}, 2023.

\bibitem[GS18]{GargS18}
Sumegha Garg and Jon Schneider.
\newblock {The Space Complexity of Mirror Games}.
\newblock In {\em Proc. 10th Conference on Innovations in Theoretical Computer
  Science}, pages 36:1--36:14, 2018.

\bibitem[HKM{\etalchar{+}}20]{HassidimKMMS20}
Avinatan Hassidim, Haim Kaplan, Yishay Mansour, Yossi Matias, and Uri Stemmer.
\newblock Adversarially robust streaming algorithms via differential privacy.
\newblock In {\em Advances in Neural Information Processing Systems 33: Annual
  Conference on Neural Information Processing Systems 2020, NeurIPS 2020,
  December 6-12, 2020, virtual}, 2020.

\bibitem[HW13]{HardtW13}
Moritz Hardt and David~P. Woodruff.
\newblock How robust are linear sketches to adaptive inputs?
\newblock In {\em Proc. 45th Annual ACM Symposium on the Theory of Computing},
  pages 121--130, 2013.

\bibitem[Ind06]{Indyk06}
Piotr Indyk.
\newblock Stable distributions, pseudorandom generators, embeddings, and data
  stream computation.
\newblock {\em J. ACM}, 53(3):307--323, 2006.

\bibitem[JP83]{JoagDevP83}
Kumar {Joag-Dev} and Frank Proschan.
\newblock Negative association of random variables, with applications.
\newblock {\em Ann. Stat.}, 11(1):286--295, 1983.

\bibitem[JW23]{JayaramW23}
Rajesh Jayaram and David~P Woodruff.
\newblock Towards optimal moment estimation in streaming and distributed
  models.
\newblock {\em ACM Trans. Alg.}, 19(3):1--35, 2023.

\bibitem[KMNS21]{KaplanMNS21}
Haim Kaplan, Yishay Mansour, Kobbi Nissim, and Uri Stemmer.
\newblock Separating adaptive streaming from oblivious streaming using the
  bounded storage model.
\newblock In {\em Advances in Cryptology - {CRYPTO} 2021 - 41st Annual
  International Cryptology Conference, {CRYPTO} 2021, Virtual Event, August
  16-20, 2021, Proceedings, Part {III}}, volume 12827 of {\em Lecture Notes in
  Computer Science}, pages 94--121. Springer, 2021.

\bibitem[MN22a]{MagenN22}
Roey Magen and Moni Naor.
\newblock Mirror games against an open book player.
\newblock In {\em 11th International Conference on Fun with Algorithms (FUN
  2022)}, volume 226, pages 20:1--20:12, 2022.

\bibitem[MN22b]{MenuhinN22}
Boaz Menuhin and Moni Naor.
\newblock Keep that card in mind: Card guessing with limited memory.
\newblock In {\em Proc. 13th Conference on Innovations in Theoretical Computer
  Science}, pages 107:1--107:28, 2022.

\bibitem[New91]{Newman91}
Ilan Newman.
\newblock Private vs.~common random bits in communication complexity.
\newblock {\em Inform. Process. Lett.}, 39(2):67--71, 1991.

\bibitem[Nis90]{Nisan90}
Noam Nisan.
\newblock Pseudorandom generators for space-bounded computation.
\newblock In {\em Proc. 22nd Annual ACM Symposium on the Theory of Computing},
  pages 204--212, 1990.

\bibitem[Nis93]{Nisan93B}
Noam Nisan.
\newblock On read once vs. multiple access to randomness in logspace.
\newblock {\em Theoretical Computer Science}, 107(1):135--144, 1993.

\bibitem[Sto23]{Stoeckl23}
Manuel Stoeckl.
\newblock Streaming algorithms for the missing item finding problem.
\newblock In {\em Proc. 34th Annual ACM-SIAM Symposium on Discrete Algorithms},
  pages 793--818, 2023.
\newblock Full version at arXiv:2211.05170v1.

\bibitem[WZ22]{WoodruffZ22}
David~P. Woodruff and Samson Zhou.
\newblock Tight bounds for adversarially robust streams and sliding windows via
  difference estimators.
\newblock In {\em Proc. 62nd Annual IEEE Symposium on Foundations of Computer
  Science}, pages 1183--1196, 2022.

\end{thebibliography}
\end{document}
